% =====================================================================
% Reviewer KuK5 -- Rating 3 (Borderline reject), Confidence 4
% Owns Rebuttal Tables 1, 2, 3 (see TABLE_REGISTRY.md).
% =====================================================================

\section*{Reviewer KuK5}

\noindent\textit{Rating 3 (Borderline reject), Confidence 4. Limitations: adequately discussed.
Ethics: none flagged. Formatting: none flagged.}

% ---------------------------------------------------------------------
\subsection*{Weaknesses}
% ---------------------------------------------------------------------

\begin{itemize}
\item[\textbf{W1.}] \itshape The neural decoding analysis cannot support a mechanistic reading
because all three causal-control protocols return null on the recovered direction and the readout
effect sizes fall in the small-to-medium band, leaving the internal evidence as correlation that
the authors themselves can only frame as a monitoring signal rather than a cause of the behavior.
\end{itemize}

\plan{기계론적 해석은 제출본이 이미 배제하고 있으므로 철회가 아니라 전면화한다. SAE 판독은 모니터링
신호로만 유지하고, 인과 통과 주장은 행동으로 정의한 축에만 한정해 읽기--쓰기 분리를 절 구조 수준에서
드러나게 다시 쓴다.}

\begin{itemize}
\item[\textbf{W2.}] \itshape The strongest behavioral claim --- that the bet-size effect is
freedom-to-choose at root rather than range expansion --- rests on a matched-cap ablation run on a
single model (GPT-4o), so its generalization to the other five depends on the broader bankruptcy
pattern rather than the same controlled test repeated across them.
\end{itemize}

\plan{6모델 매치드-캡 격자를 48/48셀 완주해 새로 실행했고, 패널 일반화 주장은 철회하되 범위확장
가설 반박은 종점이 정보를 갖는 LLaMA(캡 \$70)에 한정해 유지한다.}

% ---------------------------------------------------------------------
\subsection*{Questions}
% ---------------------------------------------------------------------

\begin{itemize}
\item[\textbf{Q1.}] \itshape Does the matched-cap dissociation (freedom-to-choose vs.\ range
expansion) hold on any of the other five models, or only GPT-4o?
\end{itemize}

\plan{Rebuttal Table 1--3으로 답하되, 사전등록 판정규칙 두 개가 모두 실패했다는 사실과 새 격자가
본문 캡 실험의 재현이 아니라는 사실을 숫자보다 먼저 밝힌다.}

\begin{itemize}
\item[\textbf{Q2.}] \itshape Given that the causal-control protocols return null, what would a
positive result on those protocols have looked like, and does their failure leave open that the
readout tracks a correlate of balance/round dynamics you did not fully residualize out?
\end{itemize}

\plan{양성 결과의 4개 기준을 명시하고, 균형·라운드 교란은 페어드 설계로 배제됨을 보이되 그 방어가
판독이 아니라 행동축에만 적용되며 중첩 기저선 검정은 미실시임을 그대로 공개한다.}

\redrule

% ---------------------------------------------------------------------
\subsection*{\color{red!60!black}Author Response}
% ---------------------------------------------------------------------

We thank you for reading both halves of the paper closely enough to locate the two places where
our evidence was thinnest: a controlled test carrying a six-model claim, and a correlational
readout at risk of being read as mechanism.

All four points ask whether our claims are broader than the tests behind them --- broader in
models than the single matched-cap ablation, and broader in mechanism than a read-side
correlation. We ran the missing grid for the first and drew the read--write boundary for the
second. Here is where they landed.

\begin{itemize}
\item \textbf{Can the neural decoding support a mechanistic reading? (W1)} No, and the submitted
paper does not make one: the abstract you read states that the neural-level analysis ``does not
claim circuit-level mechanism'', and \S4 reported both directions. The revision names the
dissociation the submitted version left implicit. What we do not withdraw is the write side ---
the behaviourally defined axis passes the pre-registered sufficiency and necessity tests, under a paired design
that holds game state fixed.
\item \textbf{Does the matched-cap claim rest on one model? (W2)} It did. We have now run the
test on all six models --- 4 caps $\times$ 2 arms at $n = 200$, all 48 cells complete --- and we
withdraw the claim that the dissociation generalises across the panel: both pre-registered
decision rules failed, and we report them as failed rather than replace them.
\item \textbf{Does the dissociation hold on any of the other five models? (Q1)} On one,
decisively. In LLaMA at cap \$70 the result strongly undercuts a pure range-expansion account: the fixed arm stakes more per
round (\$68.4 against \$32.1) yet ruins far less (3.0\% against 81.5\%). Gemini shows the same
direction too weakly to test, and the other four provide no testable contrast because neither
arm ever goes bankrupt.
\item \textbf{What would a positive causal result have looked like? (Q2)} Now stated explicitly,
as our four-criterion definition of causal control: directional change, monotone dose response,
separation from matched random directions, and stable parsing. The readout does not meet all
four. The balance/round confound cannot reach the behavioural axis, whose paired design holds
game state fixed; the nested-baseline test your wording implies has not been run, and we will
report it by 3 August.
\end{itemize}

If any part of our response falls short, we would be glad to take it further during the
discussion period.

\vspace{0.6em}

\textbf{[Weakness 2 \& Question 1]}

We understand your concern. As we understand it, you are pointing at this: our strongest
behavioural claim rests on a matched-cap ablation run on one model, so its extension to the other
five is an inference from the broader bankruptcy pattern rather than the same controlled test
repeated.

\textbf{Verdict: accepted in part.} Thank you for this sharp observation. It led us to run the
matched-cap test on all six models. We withdraw the claim that the dissociation generalises across
the panel. However, we think the range-expansion account is now undercut more sharply than before
in the one model whose endpoint is informative, LLaMA at cap \$70, and we report that evidence
below. We intend to reflect this in the camera-ready.

Two corrections come first, then the numbers.

\textbf{The first correction concerns the pre-registered rules.} Both failed, so we withdraw the
six-model generalisation claim. The primary rule is ill-posed for this design. It asks whether the
lower 2.5\% posterior quantile of the treatment coefficient exceeds zero. The fixed arm never
exceeds its variable counterpart in any cell: it reads 0.0\% in every fixed cell we tabulate
except four, and peaks at 13.0\% for LLaMA at cap \$50, against variable-arm values reaching
81.5\%. No dataset this
design produces can push that lower quantile below zero, so the rule can only pass, and passing it
is not evidence.

The registered mixed-effects model compounds the problem, because its cluster-robust standard
errors contain non-finite and divergent values. Our analysis output nevertheless marked the rule as
passing, using a pooled bootstrap interval ($+14.25$ pp $[13.25, 15.25]$) that the frozen
configuration neither specifies nor names as a fallback. That record is wrong and we withdraw it.

The qualitative secondary rule requires at least four of six models to meet two conditions at cap
\$70: wagering more than half the cap and playing more than five times as many rounds under
discretion. It returns \textbf{0 of 6}. The exposure condition alone passes in 5 of 6 models, so
the joint failure comes entirely from the stake-size condition.

A third rule, implemented in our analysis code but absent from our pre-registration, also fails, at
2 of 6. We report all three rules and do not replace them with a rule that passes. No
pre-registered rule supports a panel-level claim.

\textbf{The second correction concerns the new grid.} It does not reproduce the paper's cap
ablation, so we do not call it a reproduction. The paper's cap ablation crosses every cap with 32
prompt conditions built from five optional prompt modules, at 50 repetitions each: 128 conditions
in the variable arm and 96 in the fixed arm. The new grid keeps the \emph{base condition} alone,
the plain game prompt with none of those five modules switched on. It runs one prompt condition at
each cap, raises $n$ to 200 games per cell, and adds a \$10 fixed arm absent from the paper. Every
statement below concerns the base condition.

\textbf{Design.} The grid asks whether the matched-cap dissociation appears outside GPT-4o-mini. It
manipulates two things: the per-round cap, and whether the model sets its own stake. The
\textbf{fixed arm} sets the stake at the cap, so the model only chooses whether to play; the
\textbf{variable arm} lets it choose any stake up to the cap. We ran 6 models $\times$ 4 caps
$\times$ 2 modes at $n = 200$ games per cell, all 48 cells are complete, and a fixed-versus-variable
contrast in bankruptcy is decisive in LLaMA, weak in Gemini, and untestable in the other four.

Rebuttal Table 1 gives the bankruptcy rate of each arm at each cap, for all six models.

\vspace{0.6em}
\noindent\textbf{Rebuttal Table 1.} Matched-cap bankruptcy grid, all six models, $n = 200$ games
per cell, with 95\% Wilson intervals.

\vspace{0.3em}
\noindent
\begin{tabular}{llccc}
\toprule
\textbf{Model} & \textbf{Cap} & \textbf{Fixed \% [95\% CI]} & \textbf{Variable \% [95\% CI]} & \textbf{$\Delta$ pp} \\
\midrule
LLaMA-3.1-8B     & \$10 & 0.0 [0.0, 1.9]   & 6.5 [3.8, 10.8]   & $+6.5$ \\
                 & \$30 & 0.5 [0.1, 2.8]   & 64.0 [57.1, 70.3] & $+63.5$ \\
                 & \$50 & 13.0 [9.0, 18.4] & 71.0 [64.4, 76.8] & $+58.0$ \\
                 & \$70 & 3.0 [1.4, 6.4]   & 81.5 [75.5, 86.3] & $+78.5$ \\
\addlinespace
Gemini-2.5-Flash & \$50 & 0.5 [0.1, 2.8]   & 4.5 [2.4, 8.3]    & $+4.0$ \\
                 & \$70 & 0.0 [0.0, 1.9]   & 4.0 [2.0, 7.7]    & $+4.0$ \\
\addlinespace
GPT-4o-mini      & all  & 0.0 [0.0, 1.9]   & 0.0 [0.0, 1.9]    & 0.0 \\
GPT-4.1-mini     & all  & 0.0 [0.0, 1.9]   & 0.0 [0.0, 1.9]    & 0.0 \\
Claude-Haiku     & all  & 0.0 [0.0, 1.9]   & 0.0 [0.0, 1.9]    & 0.0 \\
Gemma-2-9B       & all  & 0.0 [0.0, 1.9]   & 0.0 [0.0, 1.9]    & 0.0 \\
\bottomrule
\end{tabular}

\vspace{0.3em}
\noindent{\small Fixed cells at caps \$30, \$50 and \$70 report post-fix values from the 18-cell
re-run after the execution defect described below. Gemini's two lower caps are not tabulated here.}
\vspace{0.6em}

LLaMA shows a large fixed-versus-variable contrast at caps \$30 to \$70, ranging from $+58.0$ to
$+78.5$ percentage points. The contrast is small at cap \$10. Gemini shows a small non-zero
contrast at the two highest caps. The other four models provide no testable contrast because
neither arm ever goes bankrupt.

\textbf{Rebuttal Table 2 gives, for each model at cap \$70, the mean stake actually executed in each
arm and the bankruptcy rate that followed. It carries the decisive result.}

Range expansion claims that discretion raises risk only because it permits larger per-round stakes.
That account predicts more ruin in the arm that stakes more per round. Our data show the opposite
ordering in the one model where both arms produce ruin. At a matched ceiling the forced arm stakes
more per round and ruins far less, so a larger available stake cannot explain the effect.

\vspace{0.6em}
\noindent\textbf{Rebuttal Table 2.} Range expansion undercut at cap \$70: the arm that stakes more
per round is the arm that ruins less.

\vspace{0.3em}
\noindent
{\small
\begin{tabular}{llccc}
\toprule
\textbf{Model} & \textbf{Arm} & \textbf{Mean executed wager} & \textbf{Wager / cap} & \textbf{Bankruptcy \% [95\% CI]} \\
\midrule
LLaMA-3.1-8B     & fixed    & \$68.4                  & ---   & 3.0 [1.4, 6.4] \\
                 & variable & \$32.1 (median \$30)    & 0.459 & 81.5 [75.5, 86.3] \\
\addlinespace
Gemini-2.5-Flash & fixed    & stake set at the \$70 cap & ---   & 0.0 [0.0, 1.9] \\
                 & variable & ---                     & 0.391 & 4.0 [2.0, 7.7] \\
\addlinespace
GPT-4o-mini      & variable & \$18.7                  & 0.267 & 0.0 [0.0, 1.9] \\
GPT-4.1-mini     & variable & \$15.1                  & 0.215 & 0.0 [0.0, 1.9] \\
Gemma-2-9B       & variable & \$18.5                  & 0.264 & 0.0 [0.0, 1.9] \\
Claude-Haiku     & variable & ---                     & 0.829 & 0.0 [0.0, 1.9] \\
\bottomrule
\end{tabular}
}

\vspace{0.3em}
\noindent{\small A dash marks a quantity we do not report for that cell. The fixed arm sets the
stake at the cap by construction. Wager / cap is the mean executed wager as a fraction of the
\$70 cap.}
\vspace{0.6em}

LLaMA's fixed arm stakes \$68.4 per round, whereas its variable arm stakes \$32.1. The fixed arm
therefore offers the larger stake. It produces 3.0\% bankruptcy, compared with 81.5\% in the
variable arm. We claim this in that model and nowhere else. Gemini's contrast points the same
way, 0.0\% fixed against 4.0\% variable, but it is too small to test the ordering, so we do not
rest the argument on it. Three further models show restraint under discretion, with mean stakes
of \$18.7, \$15.1 and \$18.5.

This restraint explains why the pre-registered stake-size clause failed in 5 of 6 models. In those
five, discretion produces wagers of 22\% to 46\% of the cap, rather than more than half. We froze
the clause before seeing any data. Its failure tests range expansion more severely than the wager
comparison alone.

\textbf{The LLaMA fixed arm is non-monotone across caps. We report the measured mechanism rather
than treating the reversal as noise.} At cap \$10, 193 of 200 games produced at least one wager,
mean rounds were 3.07, and 0 games ended bankrupt. At \$30: 170 of 200, 2.00 rounds, 1 bankrupt. At
\$50: 174 of 200, 2.07 rounds, 26 bankrupt. At \$70: 132 of 200, 0.92 rounds, 6 bankrupt.

Ruin from a \$100 balance takes two consecutive losses at cap \$50 and also two at cap \$70, so the
per-round hazard is identical at those two caps. Exposure differs. Re-betting after a first loss
falls from 57\% at cap \$30 to 35\% at cap \$50 and to 6\% at cap \$70. Withdrawal after an early
loss dominates; refusal to start does not. This weakens the simple interpretation that a larger cap
creates more risk, including the interpretation we previously used.

\textbf{The four models with no bankruptcy show a condition effect, not drift in the model versions
served by the APIs. We withdraw the drift account we gave earlier.} The paper's six-model corpus
already shows the same floor when restricted to the base condition in the variable arm: GPT-4o-mini
1/50, GPT-4.1-mini 0/50, Gemini 0/50, Claude-Haiku 0/50, Gemma 0/50, and LLaMA 22/50.

\textbf{That recomputation contradicts the scope sentence of Finding 1, and we correct Finding 1
rather than leave the two standing side by side.} Finding 1 reports a per-model bankruptcy range
shifting from $0$--$3.1\%$ under fixed betting to $5$--$72\%$ under variable, with LLaMA at
$0.4\%\!\rightarrow\!72.3\%$, and describes that comparison as made ``under the BASE prompt for all
six models''. The base-condition restriction above does not reproduce it: four models sit at 0/50,
GPT-4o-mini at 1/50, and LLaMA at 22/50, that is 44.0\% rather than 72.3\%. The quoted range is
therefore not a base-condition quantity. It must instead come from aggregation across prompt
conditions: the paper's variable-arm runner crosses 32 prompt combinations ($2^5$ over the five
optional modules named below) at every cap, and under those richer conditions variable-arm
bankruptcy reaches 100\% for Gemini (GMHW) and 78\% for Claude-Haiku (GMHWP), far above anything
the base condition produces. Those per-condition maxima are not themselves Finding 1's range, since
100\% exceeds its 72\% ceiling.

\textbf{The aggregation is the mean over those 32 conditions, and the paper prints it.} Appendix
Table~2 of the submitted paper carries the caption ``Each betting mode was aggregated from 1{,}600
games per model (32 prompt conditions $\times$ 50 repetitions)'', and its fixed-arm column reads
0.00, 0.00, 3.12, 0.00, 0.44 and 0.00 across the six models. That range, 0.00--3.12\%, is
Finding~1's ``0--3.1\%'' exactly, and the variable column reproduces its upper figure the same way.
The bankruptcy numbers are correct as computed and their provenance is already in the manuscript.

\textbf{One phrase is wrong.} Describing the comparison as made ``under the BASE prompt'' is the
error, because the quantity is the aggregate over all 32 conditions. We will restate that phrase,
together with the Figure 2a caption, in the camera-ready.

Two results close the drift account. Fisher exact tests against the new cap-\$70 variable cells are
non-significant for five of the six models ($p = 0.20$ to $1.00$). LLaMA is the only significant
difference, and it runs against drift: 44.0\% in the paper's base cells against 81.5\% now,
$p < 0.0001$, a move opposite to the floors that drift was invoked to explain. Gemma closes the
account from the other side, because it uses frozen open weights and was already at 0.0\% in the
paper's base cells, so a changed served model cannot explain its floor.

Prompt condition is what produces bankruptcy in these models. We name each condition by the letters
of the optional prompt modules it carries, drawn from the five-module set G/M/H/W/P, so that GMHWP
is the prompt carrying all five and H is the hidden-patterns module. In the paper's corpus, Gemini
rises from 0\% at base to 98\% under GMW, and GPT-4o-mini from 2\% to 44\% under GMHWP. The effect
is not uniform: LLaMA remains at 0\% under both GMHW and GMHWP.

The framing $\times$ rationality factorial reported in Reviewer gbSA's section makes the same point
at a fixed \$70 cap. A persona preamble alone moves Gemini's fixed-arm bankruptcies from 0 to 12 and
its variable-arm bankruptcies from 4 to 32. Reviewer gbSA's Weakness 2 contains our full answer,
including our retraction of the earlier claim that the preamble does essentially nothing.

\textbf{Bankruptcy is at the floor for four models, so exposure is the more informative endpoint
there.} Rebuttal Table 3 gives participation and mean rounds played in both arms at cap \$70, for
all six models. It is descriptive and was not pre-registered.

\vspace{0.6em}
\noindent\textbf{Rebuttal Table 3.} Participation and exposure at cap \$70, all six models, both
arms, $n = 200$ games per cell.

\vspace{0.3em}
\noindent
{\small
\begin{tabular}{lcccc}
\toprule
\textbf{Model} & \textbf{Fixed \%} & \textbf{Variable \%} & \textbf{Fixed rounds} & \textbf{Variable rounds} \\
\midrule
GPT-4o-mini      & 0.0\%  & 100.0\% & 0.00 & 10.28 \\
Gemini-2.5-Flash & 7.5\%  & 99.5\%  & 0.07 & 2.21 \\
LLaMA-3.1-8B     & 66.0\% & 98.0\%  & 0.92 & 15.17 \\
GPT-4.1-mini     & 0.0\%  & 89.5\%  & 0.00 & 2.05 \\
Gemma-2-9B       & 0.5\%  & 14.0\%  & 0.01 & 0.45 \\
Claude-Haiku     & 3.0\%  & 5.0\%   & 0.03 & 0.05 \\
\bottomrule
\end{tabular}
}

\vspace{0.3em}
\noindent{\small The two percentage columns report participation: the share of games with at least
one wager. The two rounds columns report mean rounds played.
\textbf{Claude-Haiku's fixed-arm figures are a parsing artefact} and do not measure behaviour. The
parser misread 5 of the 6 recorded fixed-arm wagers in this cell, and 14 of 15 across caps
\$30--\$70; all were stops. See Rebuttal Table 4 in the appendix.}
\vspace{0.6em}

\textbf{A note on naming.} We read your ``GPT-4o'' as GPT-4o-mini, the model the paper's targeted
cap ablation uses (Finding 4 and Figure 3d) and the one in our six-model corpus. In the new grid it
is one of the four floor cases, at 0.0 [0.0, 1.9] in both arms at every cap, so the original
matched-cap contrast does not reappear under the base condition at $n = 200$.

That floor is not a silent failure to replicate Finding 4, and the reconciliation lies in the
design difference stated above. The paper's cap ablation crosses each cap with all 32 prompt
conditions, so its $\sim$14/17/17\% variable-arm rates aggregate conditions in which the prompt
modules, not the base prompt, supply the bankruptcy --- in the paper's corpus GPT-4o-mini moves
from 2\% at base to 44\% under GMHWP. The new grid deliberately switches those modules off.
Restricted to the condition the two corpora share, they agree: the paper's base-condition
variable arm puts GPT-4o-mini at 1 of 50, which a Fisher exact test cannot distinguish from the
new 0 of 200. Finding 4's contrast lives in the module conditions the grid does not run, not in
a cell that failed to reproduce.

Both arms sit at the floor, so this cell distinguishes neither explanation and supports neither.
We report it because your question concerns this specific cell.

\textbf{We identified two defects in our implementation and disclose their measured impact.} The
first was a fixed-bet execution defect in the rebuttal harness. It caused fixed cells above the
lowest cap to execute the wrong stake. We found and corrected it, then re-ran those 18 cells. Every
fixed value in Rebuttal Table 1 comes from the corrected run.

The second defect affected parsing. The parser takes the first ``final decision'' match and checks
wager tokens before stop tokens, so a response that considers walking away and then stops is scored
as a wager. Re-parsing every stored response under a corrected rule flips no more than 0.3\% of the
adjudicable decisions in any corpus, and the flips run overwhelmingly from wager to stop. Rebuttal
Table 4 in the appendix gives the flip counts, the decomposed denominators and the rates for all
three corpora, and we do not repeat them here.

Two consequences belong in this section rather than in the appendix. First, 23,611 of the original
grid's 41,261 stored responses are truncated at 500 characters, so re-parsing cannot settle that
corpus. Second, run-time parsing did use the full, untruncated text; an audit of one full cell found
2 parser--text mismatches in 2,254 decisions. We read that as leaving the behavioural records
themselves unaffected.

\textbf{What we now claim, in full.} Under the base condition, discretion over stake size raises
participation in five of six models and bankruptcy in two of them. Range expansion explains neither
result. Four of the five participation increases are decisive: the fixed arm ranges from 0\% to
66\%, the variable arm from 89.5\% to 100\%. Gemma shows the fifth, from 0.5\% to 14.0\% at a mean
of 0.45 rounds, which is marginal, and we exclude Claude-Haiku because its fixed column is a parsing
artefact.

Three limits bound that claim. We do not claim that the dissociation generalises across the panel.
We do not treat the four floor cases as evidence either way. Matching the cap equalises the
per-round maximum and not the cumulative stake, so we have not separated freedom from the exposure
that freedom produces.

\textbf{We did not re-scope the goal-setting channel, and it stands as reported.} The new data do
not bear on it, and we did not re-run its matched-expected-loss control during the rebuttal.

\plan{본문의 캡 제거 실험 절은 6모델 격자로 교체하고, 표 1--3을 본문에 싣는다. 패널 일반화 문장은
삭제하고, 범위확장 반박은 LLaMA 캡 \$70에 명시적으로 한정한다. Finding 1의 ``BASE 프롬프트'' 범위
서술과 Figure 2a 캡션은 실제 집계 범위에 맞게 다시 쓰고, 이 정정을 Summary of Changes와 부록 B의
자기공개 항목에 함께 싣는다. 사전등록 규칙 3종의 실패와 고정베팅 실행 결함·파서 결함은 부록 B1·B2에
측정된 영향과 함께 공개하며, 드리프트 설명은 조건 효과 설명으로 대체한다.}

\vspace{1em}
\hrule
\vspace{1em}

\textbf{[Weakness 1 \& Question 2]}

We understand your concern. As we understand it, you are pointing at two things: the causal-control
protocols return null on the recovered direction, and the readout effect sizes are small-to-medium,
so the internal evidence remains correlational and supports monitoring rather than causation.

\textbf{Verdict: we agree about the readout direction, and the submitted paper already says so.}
Thank you for pressing on this. It led us to separate the read side from the write side at
section level and to state our four causal criteria explicitly.
The abstract you read states that the neural-level analysis ``decodes these contrasts from
decision-time internal states but does not claim circuit-level mechanism''. What that abstract did
not do is separate the two directions by name, and the revision does: it now says plainly that
SAE-decoded readout directions do not control behaviour while a behaviourally defined axis does,
and \S4's summary labels the readout the \emph{monitor} and that axis the \emph{controller}.
However, we think the premise that all causal protocols return null holds for the \emph{readout}
direction only. The paper's write tests run on a second, behaviourally defined direction, and those
are positive. Balance and round index cannot confound that result. We keep the two directions
separate throughout this reply, and we intend to reflect this separation in the camera-ready.

\textbf{The null is a read-side null. The paper's behavioural axis is positive on every write test
it reports, and this needs no new experiment.} \S4.1 reports sufficiency and necessity: steering
that axis moves Gemma's bet ratio with dose slope $z = +4.45$ against a twenty-direction chance
band, and removing it lowers betting in both models ($-0.037$ on Gemma, $p < .001$; $-0.052$ on
LLaMA, $p = .023$). \S4.2 reports cross-task write transfer on Gemma: the shared axis lowers
slot-machine betting ($z = -3.7$) while raising investment-choice risk ($z = +3.5$), each with the
sign fixed before the trials ran, for a pre-registered tally of 7 of 10 confident cells. \S4.3
reports that goal grafting raises the causal gain, at a bet-on-dose slope of $+0.0469$ against the
control's $+0.0358$.

The three protocols you count return null on the readout direction in those same tables. That
contrast is the paper's result rather than an evasion of your point: reading and writing come apart,
and only the read side is null.

\textbf{We do not withdraw a mechanistic reading of the readout, because we never made one.} The
submitted abstract disclaimed the mechanistic reading in as many words --- the neural-level
analysis ``decodes these contrasts from decision-time internal states but does not claim
circuit-level mechanism'' --- and \S4 of the submitted paper reported both directions. What it did
not do is name the dissociation, and that is what the revision adds: the readout is the monitor,
the behavioural axis is the controller. Decision-time hidden states carry a readout of the behavioural contrasts,
which supports internal-state monitoring; we did not identify a circuit, and the readout does not
show that its own direction causes the behaviour. What we change is prominence. We sharpen the
wording so the distinction is unmissable rather than carried by two sentences, and we remove the
clinical term from the title.

\textbf{We accept your characterisation of the effect sizes.} The readout effects fall in the
small-to-medium band you describe, and we do not dispute their size. We add no new effect-size
number for the readout, because we have not re-estimated it under a specification we would defend.
An effect of that size supports our retained claim, that a behavioural contrast is recoverable
above chance from decision-time states. It does not support a mechanistic claim.

\textbf{We cannot map your count of three protocols exactly onto our experiments, so we list every
protocol applied to the readout direction with its outcome.} A \emph{dose ladder} contains runs at
increasing steering strength $\alpha$. A \emph{null band} is the range of effects produced by random
directions of the same length.

\emph{(i) Steering, refit specification.} The ladder gives $+0.027$, $+0.033$ and $+0.056$ against a
null band of 0.033 from twenty matched random directions. It clears the band at one dose only.

\emph{(ii) Steering, the paper's literal Table-1 specification, ranked against eight random
directions.} $\Delta = +0.086$ at $\alpha = +2$ and $\Delta = +0.224$ at $\alpha = +4$, both with
$p < 10^{-4}$. Parse success at $\alpha = +2$ is 0.80, above our pre-set gate of 0.45, so generation
breakdown does not explain that dose. At $\alpha = +4$ parse success falls to 0.34, below the gate,
and we set that dose aside entirely. The $\alpha = +2$ cells are exploratory at $n = 50$ and are not
replicated at $n = 200$.

\emph{(iii) Removal.} Removal deletes the direction and tests whether betting falls. On the readout
direction it is null in both models, as the paper reports: $\Delta = -0.019$ with $p = 1.0$ on
LLaMA and $\Delta = -0.003$ with $p = .885$ on Gemma.

However, we think the premise that \emph{every} protocol returns null is still too strong, and the
reason is (ii) rather than (iii). Under the specification the paper itself states in Table 1,
steering the readout is significant at $\alpha = +2$ under a gate we fixed before the runs. We do
not promote that specification over the refit one, and we do not build the mechanistic reading back
on it. We also do not describe that cell as null.

\textbf{A positive result would have met four criteria, which we state as our definition of causal
control.} A direction controls betting only if all four hold. \emph{One, directional change:}
steering along the positive direction raises the betting index and along the negative direction
lowers it, with a bootstrap interval on the difference excluding zero. \emph{Two, monotone dose
response:} the effect grows across the ladder rather than appearing only at the extreme dose.
\emph{Three, separation from random directions:} the effect clears a null band built from matched
random directions of equal length. \emph{Four, stable parsing:} parse success does not degrade
differentially with dose, so generation breakdown cannot explain the effect. We fixed its threshold
before the runs at a parse-success gate of 0.45.

\textbf{The readout direction does not meet all four.} The refit ladder is monotone, so criterion
two holds. It clears the null band at one dose only, $\alpha = +3$, and in one direction, so
criterion three holds marginally at best. We have no negative-direction result, so criterion one is
untested rather than failed, and we do not count an untested criterion in our own favour. We
therefore do not classify the readout as a control lever.

\textbf{Balance and round index cannot confound the behavioural-axis result, because the design
holds game state fixed. This is the one point we do not concede.} The \textbf{behavioural axis} is
fitted from the model's own betting rather than decoded to predict it; it is not the readout. The
intervention ladders use a paired design: the same seeds run at every $\alpha$ on the same states,
and the null permutes $\alpha$ labels within each cluster. Balance and round index are therefore
identical across the compared arms, and a variable that does not differ between arms cannot generate
a difference between them.

Three further results separate the behavioural axis from a balance correlate.

\begin{itemize}
\item \textbf{The behavioural axis moves behaviour.} Its dose slope is 0.0457 at $z = +4.45$
against the twenty-direction null (null mean 0.0007, $\sigma = 0.0101$). Removing it lowers betting
by 0.037 in Gemma and by 0.052 in LLaMA, on 200 seed-matched pairs each.

\item \textbf{An axis fitted from running balance rather than betting does not.} Steering it
performs at chance ($z = +0.64$). Removing it leaves LLaMA unchanged ($\Delta = -0.006$,
$p = 1.0$) and \emph{increases} Gemma's betting ($+0.046$, $p < .001$). If balance representation
caused risk-taking, removing it would lower betting rather than raise it.

\item \textbf{The sparse autoencoder is not necessary.} A ridge probe, a regularised linear probe
fit on raw activations with no autoencoder, also steers betting, at slope 0.0284 and
$z \approx +3$. Its direction is nearly orthogonal to the readout (cos $= 0.011$), the behavioural
axis (0.021) and the balance axis (0.000).
\end{itemize}

\textbf{This evidence supports the behaviourally defined axis, not the readout.} The readout is the
direction that fails the causal test. An interpretability readout and a control lever need not
coincide, and in our data they do not.

\textbf{We have not run the nested baseline implied by your wording, and we do not claim to have
answered that part of your question.} The test would first fit a behavioural-state baseline from
balance, round index, recent outcome, choice probability and logit features, then ask whether the
hidden-state readout adds explanatory variance beyond it. We cannot yet say how much of the readout
exceeds a re-encoding of observable game state. We registered the test in Rebuttal Table 9 in the
appendix with its decision rule fixed in advance. We will report the result by 3 August, whether or
not it favours us.

Two further limits apply. Our analysis log records the LLaMA readout arm as clearing an internal
removal threshold by a margin of 6\%; the paper's significance test on that same arm returns
$\Delta = -0.019$, $p = 1.0$, so we treat that arm as null and place no weight on the margin. We
also collected post-removal projections at a layer outside the window identified by our
localisation analysis.

\textbf{We retain one deliberately limited causal claim.} We can steer and remove a
betting-aggression axis defined within this task while holding game state fixed. We did not test
causal control of loss chasing. We did not identify a circuit. None of these findings transfers to
human gambling or to deployed agents.

\plan{§3.2에서 기계론적으로 읽힐 여지가 남은 문장을 정리하고 판독을 모니터링 신호로 일관되게
재기술하며, 제목에서 임상 용어를 뺀다. 인과 통제의 4개 기준을 본문에 정의로 명시하고, 행동축(쓰기)과 판독(읽기)을 절 단위로 분리해
"세 프로토콜 null"이 판독 방향에만 해당함을 명시한다. Table-1 사양의 유의 결과와 제거 실험의 null
($\Delta=-0.019$, $p=1.0$), 분석 창 밖 레이어의 한계를 함께 공개한다. 중첩 기저선 검정은 사전
판정규칙과 함께 2차 응답에 등록한다.}
